\documentclass[11pt]{article}

\usepackage[margin=1in]{geometry}
\usepackage{amsmath,amssymb}
\usepackage{graphicx}
\usepackage{booktabs}
\usepackage{hyperref}
\usepackage{listings}
\usepackage{xcolor}
\usepackage[numbers]{natbib}

\lstset{
  basicstyle=\ttfamily\small,
  breaklines=true,
  frame=single,
  framesep=4pt,
  backgroundcolor=\color{gray!5},
  showstringspaces=false,
}

\title{Critical error rate in practice:\\context-aware variant significance for nanopore amplicon pipelines}
\author{Josh Walker}
\date{}

\begin{document}
\maketitle

\begin{abstract}
The critical error rate (CER) framework~\cite{cer_paper} provides a per-variant
significance statistic $p^*$ that is interpretable against known platform error
rates. The framework as originally posed assumes a single context-free error
rate, which is appropriate for non-homopolymer substitutions in nanopore data
but breaks down for homopolymer (HP) length variants, where empirical error
rates vary by a factor of two or more across HP length and base context~\cite{hp_paper}.
We present a context-aware extension of the framework, retaining the single
binomial equation but replacing the uniform error rate $p$ with a
context-specific $q_\text{ctx}$ derived from empirical measurement. To preserve
cross-variant comparability under context-aware $q_\text{ctx}$, we introduce
the \emph{CER factor}, $p^*/q_\text{ctx}$, as the primary derived metric. The
factor is dimensionless, comparable across variants and specimens, degrades
linearly under $q_\text{ctx}$ misspecification, and re-scales cleanly when the
empirical error table is updated. We describe the implementation in the
speconsense pipeline, including N-denominator choice for post-phasing
validation, identity-group construction, pairwise CER comparison, and a
metadata reproducibility design that allows stored output to be re-evaluated
against revised assumptions without re-running the pipeline.
\end{abstract}

\section{Introduction}

The critical error rate (CER) framework, introduced in~\cite{cer_paper}, asks
the inverse of the conventional p-value question: rather than testing whether
an observed variant is significant under an assumed error rate, it derives the
per-position error rate $p^*$ at which the observed pattern would become
plausible as artifact. A practitioner then compares $p^*$ against their
platform's known error rate; well-supported variants produce $p^*$ values
several times the platform rate, providing a transparent margin for judgment.

The framework's strength is that $p^*$ has the same units as the platform error
rate the reader already knows, decoupling the analysis from the moving target
of basecaller and chemistry improvements. A $p^*$ of $6.8\%$ is robust against
any reasonable nanopore configuration without committing to a specific assumed
rate at output time.

\subsection{The context-dependence problem}

The framework as originally formulated assumes a single per-position error rate
$p$ characterizes the platform. For non-homopolymer substitutions in current
nanopore chemistry, this is approximately true: empirical post-clustering rates
are tightly clustered around 0.6\%--1.0\% across position contexts~\cite{hp_paper}.
The framework's $p^* = 6.8\%$ headline against an assumed rate of $1.5\%$ is
robust to this spread.

For homopolymer-length variants, the assumption breaks down. Homopolymer (HP)
length errors are the dominant systematic error mode in nanopore
sequencing~\cite{delahaye2021}, and their empirical rates depend strongly on
context:

\begin{itemize}
\item HP length error rate increases monotonically from $\sim$0.7\% per position
at $L = 1$ to $\sim$1.5\% at $L = 5$, and steeply thereafter~\cite{hp_paper}.
\item G/C homopolymers consistently produce higher error rates than A/T at
matched length, with the gap widening at longer lengths.
\item Deletion-biased errors dominate insertion-biased errors at a $3$--$5\!:\!1$
ratio across moderate HP lengths.
\item Per-run and per-basecaller variation contribute additional spread of
$\sim$0.6$\times$ to $2.2\times$ on top of the context dependence.
\end{itemize}

A single assumed error rate cannot serve all of these contexts simultaneously
without either over-suppressing variants in low-error contexts (e.g., $L = 1$ HP
positions) or under-suppressing in high-error contexts (e.g., $L = 5$ G runs).
The current speconsense default of $1.5\%$ assumed error rate is calibrated to
cover the non-HP substitution regime; applied uniformly to HP variants, it
loses 30--50\% of the available signal.

\subsection{Contributions}

This paper extends the CER framework to context-dependent error rates without
abandoning its core appeal:

\begin{enumerate}
\item A \textbf{unified equation} that retains the single binomial structure
of~\cite{cer_paper} but replaces $L$ with a count of independent test sites
$n_\text{sites}$ and replaces $q$ with a context-specific $q_\text{ctx}$.
For multi-position variants ($K > 1$), the joint per-read error probability
becomes $\prod_i q_{\text{ctx},i}$ instead of $q^K$, accommodating mixed
contexts naturally.

\item A \textbf{practical decision metric}, the \emph{CER factor}, defined as
$p^* / q_\text{ctx}$. The factor is dimensionless and context-comparable by
construction, recovering the cross-variant comparability that raw $p^*$ loses
under context-aware $q_\text{ctx}$. We compare it against alternative metrics
(raw $p^*$, p-value) and recommend it as the primary annotation.

\item Implementation choices in the \textbf{speconsense pipeline} that follow
from the framework: how $N$ is chosen for post-phasing validation, how
identity groups define the comparison set for pairwise CER, why phasing
deliberately operates without statistical gating, how $q_\text{ctx}$ is
parameterized, and how HP detection is reordered relative to homopolymer
normalization.

\item A \textbf{metadata reproducibility design} that stores enough information
in each specimen's output to recompute $p^*$, $p$-value, and CER factor under
revised assumptions, without re-running the upstream pipeline.
\end{enumerate}

This paper assumes the reader is familiar with the theoretical CER framework
of~\cite{cer_paper} and the empirical HP error characterization
of~\cite{hp_paper}. Section references to these companion documents are made
where useful but the core results are not re-derived.

\section{Unified equation}
\label{sec:unified}

We retain the binomial structure of the CER paper~\cite{cer_paper}. The
generalization replaces the amplicon length $L$ with a count of independent test
sites $n_\text{sites}$ and replaces the single per-position error probability
$q$ with a context-specific $q_\text{ctx}$.

\subsection{Single-position variants ($K = 1$)}

For a candidate variant supported by $M$ reads at a single position with context
$c$ in a specimen of $N$ total reads,
\begin{equation}
\label{eq:unified_k1}
n_\text{sites} \cdot P\bigl(\mathrm{Binom}(N,\, q_{\text{ctx},c}) \geq M\bigr) = T,
\end{equation}
where $T$ is either an input significance level $\alpha$ (when solving for
$p^*$) or an output p-value (when $q_{\text{ctx},c}$ is fixed and the equation
is evaluated). This reduces to the formulation in~\cite{cer_paper} when
$q_{\text{ctx},c} = q$ is uniform across positions and
$n_\text{sites} = L$.

\subsection{Multi-position variants ($K > 1$)}

When a variant is distinguished from its comparison cluster at $K > 1$
positions, the joint per-read probability of erroring at all $K$ positions under
the null is the product of per-position rates:
\begin{equation}
\label{eq:unified_kn}
\binom{n_\text{sites}}{K} \cdot P\Bigl(\mathrm{Binom}\bigl(N,\, \textstyle\prod_{i=1}^{K} q_{\text{ctx},i}\bigr) \geq M\Bigr) = T.
\end{equation}
This generalization is essential for HP applications because mixed-context
$K > 1$ variants are common: a multi-position variant may be distinguished at
one HP position and one non-HP position simultaneously. Equation~\ref{eq:unified_kn}
reduces to $\binom{L}{K} \cdot P(\mathrm{Binom}(N, q^K) \geq M) = T$ when all
$q_{\text{ctx},i}$ are equal.

\subsection{Bonferroni site counting}
\label{sec:nsites}

The site count $n_\text{sites}$ is the number of independent positions tested.
Three definitions are operationally distinct:

\begin{description}
\item[Per-base count.] $n_\text{sites} = L$, the amplicon length in bases.
This is the formulation in~\cite{cer_paper}; appropriate when every position
is tested independently.

\item[Per-site count with HP aggregation.] $n_\text{sites} = (\text{non-HP positions}) + (\text{HP runs of } L \geq 2)$.
A length-5 HP run is one site for length-variant testing, not five.
For ITS amplicons of $\sim$700~bp, this aggregation reduces $n_\text{sites}$
modestly (typically by $\sim$15--25\%) because most HP runs are short.

\item[Per-site count with alternative-base correction.] Each non-HP position
admits 3 alternative bases; each $L \geq 2$ HP run admits multiple alternative
lengths. A strict Bonferroni would multiply $n_\text{sites}$ by the number of
plausible alternatives per site. Following~\cite{cer_paper}, we treat the $q
\to q/3$ uniform error split as implicit per-position correction over alternative
bases for non-HP positions, and use empirical per-transition rates for HP
positions (which are already specific to one alternative).
\end{description}

We adopt the per-site count with HP aggregation as the default. The choice
affects $p^*$ only logarithmically (a 25\% reduction in $n_\text{sites}$
shifts $p^*$ by less than 5\%).

\section{The CER factor}
\label{sec:factor}

\subsection{Definition}

The \emph{CER factor} is the dimensionless ratio
\begin{equation}
\label{eq:factor}
\mathrm{CER\ factor} = \frac{p^*}{q_\text{ctx}}.
\end{equation}
It answers the question: \emph{by what factor would the platform error rate at
this variant's context need to be inflated to make the observation plausible as
artifact?} A factor of 4 means quadruple the empirical error rate would be
needed; a factor of 1 corresponds exactly to the platform rate.

\subsection{Comparison with alternative metrics}

Three derived metrics summarize the same underlying CER surface:

\begin{description}
\item[Raw $p^*$.] The original CER paper's headline metric. Interpretable as
``the per-position error rate that would explain this variant.'' Comparable
across variants only when $q_\text{ctx}$ is approximately uniform.

\item[p-value.] The conventional statistic, computed by fixing $q_\text{ctx}$
and evaluating Equation~\ref{eq:unified_k1} for $T$. Universal 0--1 scale,
composable for FDR and meta-analysis, exponentially sensitive to
$q_\text{ctx}$ misspecification.

\item[CER factor.] Defined by Equation~\ref{eq:factor}. Dimensionless,
context-independent on the output side, linearly sensitive to $q_\text{ctx}$
misspecification, re-scalable when $q_\text{ctx}$ updates.
\end{description}

Table~\ref{tab:metric_comparison} summarizes the trade-offs.

\begin{table}[htbp]
\centering
\caption{Trade-off summary for derived CER metrics. ``$\checkmark$'' indicates
the metric is preferred on that dimension.}
\label{tab:metric_comparison}
\small
\begin{tabular}{lccc}
\toprule
Property & Raw $p^*$ & p-value & CER factor \\
\midrule
Context-comparable across variants                  &              &  $\checkmark$ & $\checkmark$ \\
Linear sensitivity to $q_\text{ctx}$ error          &              &              & $\checkmark$ \\
Composable for FDR / meta-analysis                  &              & $\checkmark$ &              \\
Principled default decision threshold               &              & $\checkmark$ &              \\
Reinterpretable against future chemistry            & $\checkmark$ &              & $\checkmark$ \\
Familiar to general statistical audience            &              & $\checkmark$ &              \\
Bounded dynamic range at strong end                 &              &              & $\checkmark$ \\
\bottomrule
\end{tabular}
\end{table}

\subsection{Why the factor is the right primary metric for amplicon pipelines}

For speconsense's primary use case --- amplicon barcoding where consensus
sequences are deposited into taxonomic reference databases --- the CER factor's
properties align with what practitioners need:

\begin{enumerate}
\item \textbf{Cross-variant comparability is the dominant requirement.}
Validators rank variants by confidence and apply a threshold. The factor
provides a single comparable scale across HP and non-HP contexts.

\item \textbf{Linear sensitivity to $q_\text{ctx}$ is honest.} Empirical
error rates vary $\sim$0.6$\times$ between sequencing runs on the same
basecaller~\cite{hp_paper}. P-values would swing by orders of magnitude under
this drift; the factor moves linearly, which more accurately reflects our
uncertainty about the true error rate at any given site.

\item \textbf{Re-scaling under updated chemistry is direct.} If a future
$q_\text{ctx}$ table revises the rate at a context by a factor $r$, stored
factors at that context can be rescaled by $1/r$ without recomputing from raw
$N$ and $M$. P-values must be recomputed from raw inputs; raw $p^*$ does not
require rescaling but loses cross-context comparability.

\item \textbf{The robustness interpretation survives.} A factor of 4 is
``quadruple the empirical floor would be needed to explain this as artifact''
--- a direct generalization of the raw $p^*$ interpretation, retaining the
human-interpretable units that motivated CER in the first place.
\end{enumerate}

\subsection{Default decision threshold}

We propose a default factor threshold of $\geq 2$ for accepting a variant: the
observed pattern requires at least double the empirical context-specific error
rate to be explained as artifact. The threshold is conservative under typical
$q_\text{ctx}$ uncertainty (run-to-run variation, basecaller drift) and
recovers the boundary behavior of the existing default $p^* \geq 1.5\%$
against a non-HP $q_\text{ctx} \approx 0.7\%$ (factor $\approx$ 2.1).

The threshold is a tunable: stricter calls (factor $\geq 3$) are appropriate
for novel reference deposits; permissive calls (factor $\geq 1$) are
appropriate for diversity surveys where sensitivity matters more than
specificity. Section~\ref{sec:reproducibility} describes how the factor can be
recomputed at any threshold from stored metadata without re-running the
pipeline.

\subsection{FASTA annotation scheme}
\label{sec:fasta_annotation}

Two CER fields are emitted per consensus header. The factor is the primary
decision metric; the details field is a structured string that gives the
reader enough context to interpret the factor without reproducing the full
metadata (which lives in the per-specimen JSON,
Section~\ref{sec:reproducibility}).

\begin{lstlisting}
>sample-c1 size=200 ric=100 rid=95.2 cer_factor=4.2 cer_details=p*=0.031;ctx=non-hp-sub;q=0.0059
\end{lstlisting}

For mixed-context $K > 1$ variants, the details field encodes per-position
context tags and rates joined by \texttt{+}, with the explicit $K$ value:

\begin{lstlisting}
>sample-c2 size=42 ric=42 rid=98.7 cer_factor=246.6 cer_details=p*=0.031;K=2;ctx=non-hp-sub+hp-l3-A-del1;q=0.0059+0.0213
\end{lstlisting}

The two fields:
\begin{itemize}
\item \texttt{cer\_factor}: the headline decision metric, $p^* / q_\text{ctx}$.
A single dimensionless number; higher means more robust.
\item \texttt{cer\_details}: a semicolon-separated structured string encoding
$p^*$, the $K$ value (omitted when $K = 1$), the per-position context tag(s),
and the per-position empirical error rate(s). Per-position values for $K > 1$
are joined by \texttt{+} in the order the positions were evaluated. Provides
enough information to interpret \texttt{cer\_factor} (e.g., distinguish a
single-position non-HP variant from a multi-position mixed-context variant)
and to sanity-check the factor's derivation against a remembered platform
rate. The full per-position breakdown, $\alpha$, $N$, $M$,
$n_\text{sites}$, identity-group membership, and recomputation inputs live in
the per-specimen metadata JSON (Section~\ref{sec:reproducibility}); the
header carries only what a reader needs to triage a variant at a glance.
\end{itemize}

\section{Parameterizing $q_\text{ctx}$}
\label{sec:parameterization}

The unified framework demands a context classifier and a $q_\text{ctx}$ table.
This section describes the context taxonomy, the source of empirical values,
and the operational choices for keeping $q_\text{ctx}$ current.

\subsection{Context classes}

Each MSA position falls into one of a small set of context classes. The
classification is deterministic from the consensus sequence at that position
and depends on no upstream call:

\begin{description}
\item[Non-HP substitution.] The consensus base is not part of an HP run of
length $\geq 2$, and the variant differs by a substitution. Empirical
$q_\text{ctx}$ corresponds to the per-position substitution rate; for Dorado
v5.0 on R10.4.1, $q_\text{ctx} \approx 0.0059$~\cite{hp_paper}.

\item[HP-interface substitution.] The consensus base is the first or last base
of an HP run of length $\geq 2$, and the variant is a substitution. Empirical
analysis~\cite{hp_paper}~\S 4 demonstrates that interface substitutions occur
at rates indistinguishable from non-HP substitutions; they are assigned the
non-HP substitution rate.

\item[HP length change at $(b, L, \delta)$.] The variant differs from the
consensus by a length change in an HP run of base $b$ at length $L$ by amount
$\delta$ (e.g., $\delta = -1$ for a one-base deletion). Empirical
$q_\text{ctx}$ is read from the per-(base, length, direction) table
in~\cite{hp_paper}~\S 3.3.

\item[HP length change at $L \geq 6$.] Suppressed by blanket normalization;
no $q_\text{ctx}$ assigned and no CER evaluation. This follows the
recommendation of~\cite{hp_paper}~\S 8.1: at $L \geq 6$ accuracy degrades
sharply and variant calling becomes unreliable.

\item[Indel at non-HP position.] Treated as a substitution-class event with
the non-HP rate as a placeholder. Indels at non-HP positions are rare in
amplicon data and do not warrant a separate empirical column at this stage.
\end{description}

The classifier returns a context tag and a $q_\text{ctx}$ value. For mixed-
context $K > 1$ variants, the per-position tags are recorded as a list and
the joint $q_\text{ctx}$ is the product (Equation~\ref{eq:unified_kn}).

\subsection{Sources of $q_\text{ctx}$ values}

Three operational regimes for $q_\text{ctx}$ are appropriate at different
levels of pipeline maturity:

\begin{description}
\item[Phase 1 (initial deployment): shipped per-basecaller table.] A static
$q_\text{ctx}$ table is bundled with speconsense, parameterized by basecaller
identifier (e.g., \texttt{dorado-v5.0.0}). Values are drawn from the
empirical analysis in~\cite{hp_paper} for each (base, length, direction)
combination. The table includes a version stamp recorded in metadata
(Section~\ref{sec:reproducibility}). Users select their basecaller via a CLI
flag (\texttt{--qctx-table dorado-v5.0.0}) or accept the default.

\item[Phase 2 (mature deployment): per-run estimation.] $q_\text{ctx}$ is
estimated from the current run's MSAs at pipeline startup. The estimator
pools positions across all specimens after applying the
outlier-and-bimodal-detection biological-variant filter
of~\cite{hp_paper}~\S 5.3. The estimated $q_\text{ctx}$ is recorded in
metadata so that downstream interpretation can use the run-specific values
or fall back to the shipped defaults.

\item[Override: user-supplied table.] A flag (\texttt{--qctx-table-file PATH})
accepts a TSV with columns \texttt{(context\_tag, q\_ctx)}. Useful for
research applications that need to characterize variant calls under
synthetic or extreme error assumptions.
\end{description}

The first deployment supports Phase 1 only. Phase 2 is a substantial
implementation effort and is left for future work; it is mentioned here so
that the metadata schema (Section~\ref{sec:reproducibility}) is forward-
compatible.

\subsection{Stability across runs}

The cross-run scaling factor between ont98 and run116 (same basecaller, three
months apart) is $\sim$0.6$\times$ at $L = 1$--5~\cite{hp_paper}~\S 5.4.
This means a shipped Phase 1 table may differ from a specific run's actual
rates by up to a factor of 2. The CER factor's linear sensitivity makes this
a roughly 2$\times$ shift in factor values, which is large but interpretable;
p-values under the same misspecification would shift by orders of magnitude.
Phase 2 per-run estimation eliminates this source of uncertainty.

\section{Implementation in speconsense}
\label{sec:implementation}

This section describes how the unified framework integrates with the existing
speconsense pipeline. The current implementation~\cite{speconsense} performs
recursive MSA-based variant phasing, post-phasing identity grouping, and
pairwise CER validation. The context-aware framework requires changes at each
of these stages.

\subsection{N denominator choice}
\label{sec:n_choice}

CER is evaluated only at post-phasing pairwise validation
(Section~\ref{sec:phasing_design}); phasing itself uses local frequency and
count thresholds without statistical gating. The denominator $N$ in the
binomial therefore has a single definition at the only point where it enters
the computation.

\paragraph{Identity-group sum.} When validating one cluster against another
within an identity group, $N$ is the sum of read counts within the group, not
the total specimen count. The pairwise comparison's null hypothesis is
``these two clusters represent the same underlying sequence, with the smaller
cluster arising as artifact of the larger''---a question that operates on the
population of reads sharing the same identity-group membership, not on the
broader specimen pool.

In the current implementation, this is computed in
\texttt{\_validate\_identity\_group} as
\texttt{group\_N = sum(len(subclusters[i]['read\_ids']) for i in sorted\_indices)}.

\paragraph{Why the identity-group sum is the right choice.} The artifact
hypothesis under test is ``the candidate cluster is misclassified reads from
its sibling.'' The relevant population is the union of reads that could
plausibly have ended up in either cluster --- the identity group. Using the
total specimen reads as $N$ would inflate the denominator relative to the
hypothesis under test, producing artificially small $p^*$ values (and CER
factors) and over-accepting weak variants. Using only the candidate cluster's
own size as $N$ would conversely understate the denominator and over-reject.
The identity-group sum corresponds exactly to the population over which the
candidate is hypothesized to be a redistribution of reads from validated
siblings.

\paragraph{Why not ``total specimen reads''.} A natural alternative is to use
the total specimen input reads as $N$, on the conservative reasoning that any
read in the specimen could in principle have appeared in this position. We
considered and rejected this. The artifact null is not ``any specimen read
could have erred here'' --- which is what total specimen $N$ models --- but
``among the reads belonging to this identity group, could redistribution have
produced the candidate's count.'' Reads outside the group are part of
unrelated identity groups and do not enter the artifact mechanism for this
candidate.

\subsection{Identity groups}
\label{sec:identity_groups}

An \emph{identity group} is a connected component of clusters whose pairwise
HP-normalized identity exceeds a threshold. Groups define the comparison set
for pairwise CER: a candidate variant is tested against members of its own
identity group, not against the entire specimen.

\paragraph{Construction.} Groups are formed by union-find using
\texttt{\_hp\_normalized\_pairwise\_compare} as the distance function. The
default threshold is 85\% identity (15\% adjusted-identity distance), tunable
via \texttt{--group-identity}. Two clusters are joined if their pairwise
identity exceeds the threshold; transitive closure forms groups.

\paragraph{Role in CER decisions.} Within each identity group, clusters are
ordered by size. The largest cluster (the \emph{anchor}) auto-passes CER
without comparison. Each smaller cluster is tested pairwise against all
already-validated clusters in the group; the candidate's CER factor is the
\emph{minimum} across all pairwise comparisons, corresponding to the
``nearest plausible artifact source.'' Acceptance requires the minimum
factor to exceed the decision threshold.

\paragraph{Why groups, not any-to-any.} The artifact hypothesis is that a
small cluster represents reads misclassified from a similar cluster. This
hypothesis only makes sense within a similarity group; clusters of distinct
sequences cannot serve as artifact sources for each other. Restricting the
comparison set to identity-group members both encodes this hypothesis and
reduces compute by avoiding $O(n^2)$ comparisons across the specimen.

\subsection{Pairwise CER comparison}
\label{sec:pairwise}

Pairwise CER between candidate cluster $A$ and validated reference cluster $B$
proceeds as follows:

\begin{enumerate}
\item Compute pairwise differences $K = \mathrm{count\_variant\_differences}(A_\text{cons}, B_\text{cons})$,
where each substitution contributes 1 and each contiguous indel event
contributes 1 (regardless of indel length).

\item For each of the $K$ differing positions, classify the context and look
up $q_{\text{ctx},i}$.

\item Compute the joint per-read error probability $\prod_i q_{\text{ctx},i}$.

\item Take $L = \max(|A_\text{cons}|, |B_\text{cons}|)$ as the consensus
length input to $n_\text{sites}$ (which subtracts HP-aggregated runs as
described in Section~\ref{sec:nsites}).

\item Take $N$ as the identity-group sum (Section~\ref{sec:n_choice}) and $M$
as the candidate cluster's read count.

\item Solve Equation~\ref{eq:unified_kn} for $p^*$ at the configured $\alpha$,
then compute the CER factor.
\end{enumerate}

The candidate's reported CER factor is the minimum over all pairwise
comparisons against validated cluster references. This conservative choice
reflects the worst-case artifact source and ensures that if any nearby
validated cluster could plausibly explain the candidate as artifact, the
candidate is suppressed.

\subsection{Phasing without statistical gating}
\label{sec:phasing_design}

Recursive MSA-based phasing in speconsense produces candidate clusters using
only local frequency and count thresholds (\texttt{--min-variant-frequency},
\texttt{--min-variant-count}); CER is not evaluated during phasing. This is
a deliberate design choice, not an oversight, and we intend to keep it that
way under the context-aware framework.

\paragraph{Phasing as candidate generation.} The role of phasing is to
\emph{partition} reads into candidate haplotype clusters based on positional
allele patterns. Phasing answers the structural question ``which reads
co-vary at which positions,'' producing a set of candidate clusters whose
existence is a hypothesis to be tested downstream. The role of CER is to
\emph{evaluate} those candidates against the artifact null. Conflating
generation and evaluation forces every phasing decision through the CER
machinery, which is both expensive and structurally inappropriate: at the
moment a phasing split is being considered, the candidate cluster has not
yet had its consensus generated, so its $K$ relative to siblings is unknown
and the appropriate $q_\text{ctx}$ values cannot be computed.

\paragraph{Why post-phasing CER is preferred.} Three reasons:

\begin{enumerate}
\item \textbf{Mixed-context $K > 1$ requires consensus alignment.} The
unified equation requires per-position $q_{\text{ctx},i}$ values for each of
the $K$ positions distinguishing the candidate from its sibling. These are
known only after both consensuses are generated and aligned, which happens
post-phasing. A phasing-time gate would be limited to $K = 1$ evaluation
against an arbitrary reference, losing the substantial sensitivity gains of
multi-position evaluation (Section~\ref{sec:unified}).

\item \textbf{Phasing's local thresholds are fast and effective pre-filters.}
The existing \texttt{--min-variant-count 5} and
\texttt{--min-variant-frequency 0.10} thresholds eliminate the
overwhelmingly majority of trivial splits (a single-read singleton at a
random position) without any statistical machinery. CER's contribution is
discriminating among the remaining candidates, all of which have at least
modest local support. Running CER at every phasing step would be largely
wasted work.

\item \textbf{Identity-group $N$ depends on phasing's output.} The
identity-group sum used as $N$ in pairwise CER (Section~\ref{sec:n_choice})
is determined by the set of clusters phasing produces and how they group by
similarity. The hypothesis a phasing-time gate would test (``is this
sub-split significant against \emph{some} pool'') has no well-defined $N$
until the full set of clusters is known.

\item \textbf{Cleaner primary consensuses via noise excision.} Indiscriminate
phasing followed by post-validation rejection lets non-significant variants
be cleanly identified as their own \emph{ns} (not significant) clusters
rather than being absorbed back into a parent and manifesting as IUPAC
ambiguity codes. With a phasing-time gate, reads that would have formed a
non-significant split remain in their parent cluster, contributing to
ambiguity at the affected positions and dirtying the parent's consensus. The
no-gate design separates ``structurally a candidate variant'' from
``statistically meaningful,'' so a primary consensus accumulates ambiguity
only from positions where reads genuinely could not be phased apart at
all---not from positions where phasing succeeded but the resulting cluster
was too weakly supported to report. The non-significant reads remain visible
in the output as ns clusters, available for inspection without polluting the
significant consensuses.
\end{enumerate}

\paragraph{What about over-splitting?} A reasonable concern is that
indiscriminate phasing produces clusters that should never have been split,
inflating downstream compute and adding noise. In practice, the local
frequency and count thresholds do most of this work: a 10\% / 5-read floor
prevents single-read splits and cleanly catches obvious noise. Post-phasing
identity-group construction (Section~\ref{sec:identity_groups}) collapses
near-identical splits via union-find, and pairwise CER then suppresses
weakly-supported candidates. The pipeline's three-stage structure
(local-threshold phasing $\to$ similarity-based grouping $\to$ statistical
validation) cleanly separates concerns and matches each stage to the
question it is structurally suited to answer.

\subsection{HP detection and normalization}
\label{sec:hp_normalization}

The existing pipeline applies HP-length normalization in two places:

\begin{enumerate}
\item During clustering identity computation (\texttt{\_hp\_normalized\_pairwise\_compare}),
HP length differences within runs of length $\geq 6$ contribute zero distance.
\item During MSA-based variant detection, HP length differences are blanket-
discarded as ``not real variants'' regardless of length.
\end{enumerate}

The first behavior is preserved: HP normalization in clustering distance is
a separate concern from variant calling and need not change.

The second behavior changes when the context-aware framework is active. HP
length variants at $L \leq 5$ become candidate variants subject to CER
evaluation with the appropriate $q_\text{ctx}$. HP length variants at
$L \geq 6$ continue to be blanket-discarded, following the recommendation
of~\cite{hp_paper}~\S 8.1.

\paragraph{Pipeline ordering.} The MSA variant detector must be reordered
so that HP length differences are surfaced as candidate variants (with their
context tags) rather than dropped. The CER validation phase then evaluates
these candidates with the same machinery as substitution variants.
Implementation note: the existing
\texttt{is\_variant\_position\_with\_composition} function tracks
\texttt{homopolymer\_composition} separately from raw base counts; this
information feeds the context classifier directly.

\paragraph{Ambiguity handling.} HP length variants that fail CER (factor $< 2$)
should be absorbed into the parent cluster's consensus as ambiguity, mirroring
the behavior for failed substitution variants. Whether to encode HP length
ambiguity as IUPAC codes, as the modal length, or as a special annotation is
an open question; the current implementation uses the modal length without
explicit ambiguity annotation. A revised approach could record HP length
variation (e.g., \texttt{poly\_A[5..6]} as a header annotation) for downstream
inspection.

\section{Metadata reproducibility}
\label{sec:reproducibility}

The CER computation depends on several inputs --- $N$, $M$, $K$,
$n_\text{sites}$, $q_\text{ctx}$, $\alpha$ --- and a stable mapping from
position to context. As empirical $q_\text{ctx}$ values are revised (new
basecaller versions, per-run estimation, improved error characterization),
practitioners should be able to recompute CER factor and p-value from
\emph{stored} metadata, without re-running the pipeline. This section
specifies the metadata schema that supports this.

\subsection{Specimen metadata JSON}

Each specimen produces a metadata JSON file at the standard output location.
The current schema~\cite{speconsense} contains parameters, input information,
and primer settings; we extend it with CER reproduction data:

\begin{lstlisting}
{
  "schema_version": "2.0",
  "sample_name": "BCFCQ-25-005",
  "timestamp": "2026-04-16T12:34:56Z",
  "parameters": {
    ...existing fields...,
    "assumed_error_rate": 0.015,
    "significance_level": 1e-5,
    "qctx_table": "dorado-v5.0.0",
    "qctx_table_version": "20251201",
    "qctx_source": "shipped"
  },
  "specimen": {
    "total_input_reads": 1000,
    "consensus_length": 698,
    "n_sites": 587,
    "n_sites_breakdown": {
      "non_hp": 423,
      "hp_runs": 164
    }
  },
  "identity_groups": [
    {
      "group_id": "g1",
      "group_N": 942,
      "members": ["c1", "c2", "c3"]
    }
  ],
  "variants": [
    {
      "cluster_id": "c2",
      "identity_group": "g1",
      "anchor": false,
      "M": 42,
      "K": 2,
      "L": 698,
      "context_tags": ["non-hp-sub", "hp-l3-A-del1"],
      "q_ctx_per_position": [0.0059, 0.0213],
      "q_ctx_joint": 1.257e-4,
      "compared_against": "c1",
      "cer_pstar": 0.031,
      "cer_factor": 5.3,
      "p_value_at_qctx": 2.7e-9
    }
  ]
}
\end{lstlisting}

\subsection{Recomputation workflow}

Given a stored metadata JSON and a revised $q_\text{ctx}$ table or significance
level, recomputation proceeds:

\begin{enumerate}
\item For each variant, look up the new $q_{\text{ctx},i}$ values for the
recorded \texttt{context\_tags}.
\item Compute the new joint $q_\text{ctx,joint} = \prod_i q_{\text{ctx},i}^\text{new}$.
\item Recompute $p^*$ from $(N, M, n_\text{sites}, K, \alpha^\text{new})$.
\item The CER factor becomes $p^* / q_\text{ctx,joint}^\text{new}$ where
$q_\text{ctx,joint}^\text{new}$ is the geometric mean of the new per-position
rates raised to $1/K$, or equivalently a per-position factor based on the
arithmetic structure preferred for the application.
\item P-value is $n_\text{sites} \cdot P(\mathrm{Binom}(N, q_\text{ctx,joint}^\text{new}) \geq M)$.
\end{enumerate}

The original $N$, $M$, and consensus-derived $n_\text{sites}$ are immutable
properties of the sequencing data; no upstream pipeline run is needed to
revise the significance metric.

\subsection{Schema versioning}

The \texttt{schema\_version} field allows downstream tools to detect missing
fields and degrade gracefully. Version 1.0 is the pre-context-aware schema;
version 2.0 introduces the CER fields in this paper. Tools reading version 1.0
output can fall back to the original raw-$p^*$ interpretation; tools reading
version 2.0 output have full reproducibility.

\subsection{Why this is worth the bytes}

The metadata addition adds $\sim$200--500 bytes per variant. For a typical
ITS amplicon study with 1000 specimens and an average of 3 variants each, the
total overhead is on the order of 1--2~MB across the run. This is negligible
relative to the FASTA and FASTQ output. The benefit --- recomputation under
revised assumptions without re-running --- is substantial: as basecaller
versions evolve and empirical error tables improve, every stored output
remains scientifically valid without recomputation cost.

\section{Validation}
\label{sec:validation}

\textbf{This section is a placeholder. Empirical content will be added once
Phase 1 implementation generates data.}

The planned validation reproduces the structure of the original CER
validation~\cite{cer_paper}~\S 8 but with the context-aware framework:

\begin{itemize}
\item Apply context-aware CER to the ONT98 dataset (837 specimens) and
compare per-variant decisions against the existing raw-$p^*$ implementation
at default settings.

\item Apply to run116 (127 specimens, herbarium type material) and quantify
the impact of the cross-run $q_\text{ctx}$ scaling factor identified
in~\cite{hp_paper}~\S 5.4 on accept/reject decisions.

\item Sensitivity analysis: vary $q_\text{ctx}$ by $\pm 50\%$ at each context
class and quantify the fraction of decisions that flip. This characterizes the
framework's robustness to $q_\text{ctx}$ misspecification under realistic
uncertainty.

\item HP-variant rescue: identify HP length variants that the current pipeline
discards via blanket normalization but the context-aware framework would
accept. Manual review of a sample to assess biological plausibility.

\item Specimen-level outcomes: confirm that no specimens are lost (zero
organism dropouts), consistent with the original CER validation.
\end{itemize}

Expected results: a moderate increase in HP-variant calls (mostly $L = 1$--$3$
variants currently discarded), a modest reduction in marginal substitution
variants at HP-interface positions (where the empirical $q_\text{ctx}$ is
slightly lower than the current uniform assumption), and stable specimen-level
counts.

\section{Limitations}

The framework as presented inherits the limitations of the underlying CER
paper~\cite{cer_paper} and adds a few specific to the context-aware extension.

\paragraph{Cross-read error correlation.} The binomial null assumes
independence of errors across reads. Reads sharing the same PCR template can
exhibit correlated errors (PCR-introduced substitutions, amplification bias);
these are physically present in the library and not addressed by the CER
framework, which establishes physical library presence rather than biological
reality.

\paragraph{Within-read non-independence at HP positions.} HP length errors are
signal-level events affecting the entire run at once, not independent per-base
errors~\cite{hp_paper}~\S 3.4. The per-position $q_\text{ctx}$ tabulated in
the empirical analysis is a compact parameterization for cross-context
comparison; the binomial model used here is correct because it operates
across reads (independent draws from the sequencing process), not within a
read. Practitioners should not interpret the per-position $q_\text{ctx}$ for
HP positions as a mechanistic rate.

\paragraph{Non-stationary error rates within a run.} If error rates drift
during a sequencing run (flow cell degradation, temperature variation), the
binomial model's identical-trials assumption is violated. Phase 2 per-run
estimation partially mitigates this for between-run drift but cannot address
within-run drift.

\paragraph{$q_\text{ctx}$ provenance for novel basecaller versions.} Shipped
$q_\text{ctx}$ tables are derived from empirical analysis of specific
basecaller versions~\cite{hp_paper}. New basecaller releases require either
re-characterization or use of the closest available table with explicit
acknowledgment of the mismatch. Per-run estimation (Phase 2) addresses this
but is not yet implemented.

\paragraph{Long HP runs ($L \geq 6$).} The framework does not extend to
$L \geq 6$ HP variants; these continue to be blanket-discarded as in the
existing pipeline. The HP error analysis~\cite{hp_paper}~\S 8.1 documents that
accuracy degrades sharply at $L \geq 6$ and variant calling becomes
unreliable; this is a property of the data, not the framework.

\paragraph{Biological variant priors.} The framework treats all positions as
having the same prior probability of carrying a real biological variant. In
ITS, conserved regions (5.8S) and hypervariable regions (ITS1, ITS2) have
very different prior variant densities. A Bayesian extension could incorporate
position-specific priors; the current framework does not.

\section{Conclusion}

The CER framework~\cite{cer_paper} provides a transparent statistic for variant
significance that is interpretable against platform error knowledge. Its
extension to context-dependent error rates --- demanded by the empirical
characterization of HP errors~\cite{hp_paper} --- requires changes to the
inputs of the CER equation, not to its structure. The unified equation
(Section~\ref{sec:unified}) accommodates per-context error rates and mixed-
context multi-position variants without abandoning the binomial framework.

The CER factor (Section~\ref{sec:factor}) is the practical metric that recovers
cross-variant comparability under context-aware $q_\text{ctx}$. It preserves
the human-interpretable units that motivated the original framework, degrades
linearly under realistic $q_\text{ctx}$ uncertainty, and re-scales cleanly
when the empirical error table is updated.

The implementation in speconsense (Section~\ref{sec:implementation}) integrates
the framework with existing pipeline structure: the $N$ denominator follows
the conservative choice at each stage, identity groups define the pairwise
comparison set, and HP variant detection is reordered so that $L \leq 5$ HP
length differences become candidate variants subject to CER rather than being
unconditionally discarded.

The metadata reproducibility design (Section~\ref{sec:reproducibility}) ensures
that stored output can be reinterpreted under updated assumptions. As
basecaller chemistry evolves and empirical error tables improve, no
recomputation from raw reads is required --- the pipeline's conclusions remain
adaptable to the moving target of nanopore error characterization.

Phase 2 work --- per-run $q_\text{ctx}$ estimation, an empirical-distribution
model for HP variants that captures within-read signal-level dependencies, and
incorporation of biological-variant priors --- remains future work, scoped
beyond the initial context-aware deployment.

\subsection*{Acknowledgments}

Substantial portions of the text were drafted with the assistance of Claude
(Anthropic, Opus 4.7). The core ideas, framework design, and implementation
choices were provided by the author; the AI contributed to exposition,
\LaTeX\ formatting, and iterative drafting under human direction. The author
takes full responsibility for the content.

\appendix

\section{Metadata JSON schema (full example)}
\label{app:schema}

A minimal but complete example showing all CER-related fields. Comments
(prefixed with \texttt{//}) are illustrative and not part of the JSON.

\begin{lstlisting}
{
  "schema_version": "2.0",
  "sample_name": "BCFCQ-25-005",
  "timestamp": "2026-04-16T12:34:56Z",
  "speconsense_version": "0.7.0",

  "parameters": {
    "algorithm": "mcl",
    "min_identity": 0.90,
    "presample_size": 1000,
    "assumed_error_rate": 0.015,
    "significance_level": 1e-5,
    "factor_threshold": 2.0,
    "qctx_table": "dorado-v5.0.0",
    "qctx_table_version": "20251201",
    "qctx_source": "shipped"
  },

  "input_file": "BCFCQ-25-005.fastq.gz",
  "specimen": {
    "total_input_reads": 1000,
    "consensus_length": 698,
    "n_sites": 587,
    "n_sites_breakdown": {
      "non_hp_positions": 423,
      "hp_runs_l2_plus": 164
    }
  },

  "identity_groups": [
    {
      "group_id": "g1",
      "group_N": 942,
      "members": ["c1", "c2", "c3"]
    },
    {
      "group_id": "g2",
      "group_N": 58,
      "members": ["c4"]
    }
  ],

  "variants": [
    {
      "cluster_id": "c1",
      "identity_group": "g1",
      "anchor": true,
      "size": 850,
      "ric": 100,
      "M": 850,
      "cer_pstar": null,
      "cer_factor": null,
      "note": "anchor: auto-passes CER"
    },
    {
      "cluster_id": "c2",
      "identity_group": "g1",
      "anchor": false,
      "size": 50,
      "ric": 50,
      "M": 50,
      "K": 1,
      "L": 698,
      "context_tags": ["non-hp-sub"],
      "q_ctx_per_position": [0.0059],
      "q_ctx_joint": 0.0059,
      "compared_against": "c1",
      "cer_pstar": 0.087,
      "cer_factor": 14.7,
      "p_value_at_qctx": 1.2e-15,
      "decision": "accept"
    },
    {
      "cluster_id": "c3",
      "identity_group": "g1",
      "anchor": false,
      "size": 42,
      "ric": 42,
      "M": 42,
      "K": 2,
      "L": 698,
      "context_tags": ["non-hp-sub", "hp-l3-A-del1"],
      "q_ctx_per_position": [0.0059, 0.0213],
      "q_ctx_joint": 1.257e-4,
      "compared_against": "c1",
      "cer_pstar": 0.031,
      "cer_factor": 246.6,
      "p_value_at_qctx": 4.1e-9,
      "decision": "accept"
    }
  ]
}
\end{lstlisting}

\section{Default $q_\text{ctx}$ table for Dorado v5.0}
\label{app:qctx_table}

Drawn from~\cite{hp_paper} Tables 1, 2, and 4. Values are filtered per-position
error rates (after biological-variant filtering) for the ont98 dataset
(R10.4.1, Dorado SUP v5.0.0). Per-(base, length, direction) values for HP
length changes are derived from the supplementary distributions; the table
below shows summary values.

\begin{table}[htbp]
\centering
\caption{Default $q_\text{ctx}$ values for Dorado v5.0 contexts (per-position
total error rate). Context tag convention: \texttt{hp-l\{L\}-\{base\}} for HP
length changes; direction-specific subtags (\texttt{-del1}, \texttt{-ins1})
override the aggregate when available.}
\label{tab:default_qctx}
\small
\begin{tabular}{lcl}
\toprule
Context tag & $q_\text{ctx}$ & Source \\
\midrule
\texttt{non-hp-sub}            & 0.0059 & \cite{hp_paper} Table 5 (ont98) \\
\texttt{non-hp-indel}          & 0.0108 & sum of Del + Ins, \cite{hp_paper} Table 5 \\
\texttt{hp-interface-sub}      & 0.0059 & treated as non-HP per \cite{hp_paper}~\S 4.3 \\
\midrule
\texttt{hp-l1-A}               & 0.0067 & \cite{hp_paper} Table 4 \\
\texttt{hp-l1-T}               & 0.0067 & \cite{hp_paper} Table 4 (pooled) \\
\texttt{hp-l1-C}               & 0.0067 & \cite{hp_paper} Table 4 (pooled) \\
\texttt{hp-l1-G}               & 0.0067 & \cite{hp_paper} Table 4 (pooled) \\
\texttt{hp-l2-A}               & 0.0083 & \cite{hp_paper} Table 4 \\
\texttt{hp-l3-A}               & 0.0097 & \cite{hp_paper} Table 4 \\
\texttt{hp-l4-A}               & 0.0099 & \cite{hp_paper} Table 4 \\
\texttt{hp-l5-A}               & 0.0113 & \cite{hp_paper} Table 4 \\
\texttt{hp-l5-G}               & 0.0140 & per-base, \cite{hp_paper} Table 10 \\
\midrule
\texttt{hp-l6+}                & N/A & blanket discard, \cite{hp_paper}~\S 8.1 \\
\bottomrule
\end{tabular}
\end{table}

The shipped Phase 1 table will include direction-specific subtags
(\texttt{-del1}, \texttt{-ins1}, \texttt{-del2+}) and per-base values at all
lengths; Table~\ref{tab:default_qctx} shows summary values for illustration.
The precise table is generated from the supplementary
\texttt{hp\_distributions.tsv} files of~\cite{hp_paper} via a build script
that ships with speconsense.

\begin{thebibliography}{99}

\bibitem[Walker, 2026a]{cer_paper}
J.~Walker.
\newblock Critical error rate analysis for {MSA}-based variant phasing in
nanopore amplicon sequencing.
\newblock Technical report, 2026.

\bibitem[Walker, 2026b]{hp_paper}
J.~Walker.
\newblock Homopolymer error rate analysis for nanopore amplicon sequencing.
\newblock Technical report, 2026.

\bibitem[Walker, 2024]{speconsense}
J.~Walker.
\newblock Speconsense: High-quality clustering and consensus generation for
{Oxford Nanopore} amplicon reads.
\newblock \url{https://github.com/joshuaowalker/speconsense}, 2024.

\bibitem[Delahaye and Nicolas, 2021]{delahaye2021}
C.~Delahaye and J.~Nicolas.
\newblock Sequencing {DNA} with nanopores: Troubles and biases.
\newblock \emph{PLoS ONE}, 16(10):e0257521, 2021.

\bibitem[Sereika et~al., 2022]{sereika2022}
M.~Sereika, R.~H. Kirkegaard, S.~M. Karst, T.~Y. Michaelsen, E.~A.
S{\o}rensen, R.~D. Wollenberg, and M.~Albertsen.
\newblock {Oxford Nanopore R10.4} long-read sequencing enables the generation of
near-finished bacterial genomes from pure cultures and metagenomes without
short-read or reference polishing.
\newblock \emph{Nature Methods}, 19:482--488, 2022.

\bibitem[Wick et~al., 2023]{wick2023}
R.~R. Wick, L.~M. Judd, and K.~E. Holt.
\newblock Assembling the perfect bacterial genome using {Oxford Nanopore} and
{Illumina} sequencing.
\newblock \emph{PLOS Computational Biology}, 19(3):e1010905, 2023.

\bibitem[Lee, 2003]{lee2003}
C.~Lee.
\newblock Generating consensus sequences from partial order multiple sequence
alignment graphs.
\newblock \emph{Bioinformatics}, 19(8):999--1008, 2003.

\bibitem[Vaser et~al., 2017]{vaser2017}
R.~Vaser, I.~Sovi\'{c}, N.~Nagarajan, and M.~\v{S}iki\'{c}.
\newblock Fast and accurate de novo genome assembly from long uncorrected reads.
\newblock \emph{Genome Research}, 27(5):737--746, 2017.

\end{thebibliography}

\end{document}
